%%%%%%%%%%%%%%%%%%%%%%%%%%%%%%%%%%%%%%%%%%%%%%%%%%%%%%%%%%%%%%%%%%%%%%%%
%% DADOS DO PEDIDO — o único arquivo de metadados do template.         %%
%%                                                                     %%
%% Lido pelas quatro raízes (relatorio-descritivo.tex,                 %%
%% reivindicacoes.tex, desenhos.tex e resumo.tex), de modo que o        %%
%% título seja necessariamente o mesmo no relatório descritivo e no     %%
%% resumo — como exige o art. 24, IV da Portaria/INPI/DIRPA nº 14/2024. %%
%%%%%%%%%%%%%%%%%%%%%%%%%%%%%%%%%%%%%%%%%%%%%%%%%%%%%%%%%%%%%%%%%%%%%%%%

%%%%%%%%%%%%%%%%%%%%%%%%%%%%%%%%%%%%%%%%%%%%%%%%%%%%%%%%%%%%%%%%%%%%%%%%
%% 1. TÍTULO DO PEDIDO
%%
%% Art. 24 — o título deve:
%%   I   ser conciso, claro e específico, identificando o objeto do pedido,
%%       não podendo exceder 500 caracteres;
%%   II  não conter denominações de fantasia;
%%   III não conter fórmulas químicas ou matemáticas;
%%   IV  ser o mesmo no relatório descritivo e no resumo.
%% Art. 25 — deve representar adequadamente as diferentes categorias de
%%       reivindicação (se você reivindica produto E processo, o título
%%       precisa refletir os dois).
%%%%%%%%%%%%%%%%%%%%%%%%%%%%%%%%%%%%%%%%%%%%%%%%%%%%%%%%%%%%%%%%%%%%%%%%

% Os \providecommand permitem que o showcase (./gerar-exemplos.sh) injete
% outros valores por linha de comando, sem duplicar nenhum arquivo. Na
% compilação normal valem os valores definidos aqui — é aqui que você edita.
\providecommand{\TituloDoPedido}{DISPOSITIVO DE ACIONAMENTO PARA VÁLVULA DE
IRRIGAÇÃO E PROCESSO DE CONTROLE DE VAZÃO}
\titulodopedido{\TituloDoPedido}

%%%%%%%%%%%%%%%%%%%%%%%%%%%%%%%%%%%%%%%%%%%%%%%%%%%%%%%%%%%%%%%%%%%%%%%%
%% 2. NATUREZA DA PROTEÇÃO
%%
%%   invencao             — patente de invenção. Reivindicações conforme os
%%                          arts. 29 a 31; desenhos opcionais.
%%   modelo-de-utilidade  — patente de modelo de utilidade. UMA ÚNICA
%%                          reivindicação independente (art. 33), dependentes
%%                          só nas hipóteses do art. 35, e desenhos
%%                          OBRIGATÓRIOS (art. 22).
%%
%% Se você não sabe qual é o seu caso: invenção é uma solução técnica nova
%% para um problema técnico (art. 8º da LPI); modelo de utilidade é um objeto
%% de uso prático, ou parte dele, com nova forma ou disposição que resulte em
%% melhoria funcional no seu uso ou fabricação (art. 9º da LPI).
%%%%%%%%%%%%%%%%%%%%%%%%%%%%%%%%%%%%%%%%%%%%%%%%%%%%%%%%%%%%%%%%%%%%%%%%

\providecommand{\NaturezaDoPedido}{invencao}
\natureza{\NaturezaDoPedido}

%%%%%%%%%%%%%%%%%%%%%%%%%%%%%%%%%%%%%%%%%%%%%%%%%%%%%%%%%%%%%%%%%%%%%%%%
%% 3. MODALIDADE DO DEPÓSITO
%%
%%   originario         — depósito comum. Nada é impresso após o título.
%%   adicao             — certificado de adição de invenção. Imprime, após o
%%                        título, "Certificado de adição de invenção do ___"
%%                        (art. 43, II). Só existe para natureza 'invencao'
%%                        (arts. 42 e 43) e exige \pedidovinculado.
%%   divisao            — pedido dividido. Imprime, após o título,
%%                        "Dividido do ___" (art. 51, II). Exige
%%                        \pedidovinculado.
%%   fase-nacional-pct  — entrada na fase nacional de pedido internacional
%%                        (Portaria INPI nº 39/2021). A forma dos documentos é
%%                        a mesma desta norma (art. 65); os dados do pedido
%%                        internacional vão no requerimento eletrônico, não
%%                        nos documentos.
%%%%%%%%%%%%%%%%%%%%%%%%%%%%%%%%%%%%%%%%%%%%%%%%%%%%%%%%%%%%%%%%%%%%%%%%

\providecommand{\ModalidadeDoPedido}{originario}
\modalidade{\ModalidadeDoPedido}

% Número do pedido ou da patente a que este documento se vincula. Obrigatório
% nas modalidades 'adicao' e 'divisao'; ignorado nas demais.
% Exemplo: \pedidovinculado{BR 10 2024 000000 0}
\providecommand{\PedidoVinculado}{}
\pedidovinculado{\PedidoVinculado}

%%%%%%%%%%%%%%%%%%%%%%%%%%%%%%%%%%%%%%%%%%%%%%%%%%%%%%%%%%%%%%%%%%%%%%%%
%% 4. FORMATO DO RÓTULO DOS PARÁGRAFOS
%%
%% Art. 26, II — cada parágrafo do relatório descritivo recebe numeração
%% sequencial em algarismos arábicos à esquerda do texto. A norma exemplifica
%% com três dígitos ([003], [015]); os formulários oficiais do e-Patentes usam
%% um dígito ([1], [2]). As duas formas atendem à norma.
%%
%%   tres-digitos  — [001], [002], ... (padrão)
%%   arabico       — [1], [2], ...
%%%%%%%%%%%%%%%%%%%%%%%%%%%%%%%%%%%%%%%%%%%%%%%%%%%%%%%%%%%%%%%%%%%%%%%%

\formatorotuloparagrafo{tres-digitos}

%%%%%%%%%%%%%%%%%%%%%%%%%%%%%%%%%%%%%%%%%%%%%%%%%%%%%%%%%%%%%%%%%%%%%%%%
%% 5. CÓPIA DE COMPARAÇÃO (arts. 51, III e 57, II)
%%
%% Ao apresentar modificação depois do depósito, ou ao depositar pedido
%% dividido, a norma pede dois documentos tirados da MESMA matéria: os
%% documentos modificados "sem qualquer tipo de rasura ou sinalização"
%% (art. 57, I) e uma cópia de comparação com tachado indicando remoção e
%% sublinhado indicando inclusão ou substituição (art. 57, II).
%%
%% Marque as alterações no texto com \removido{...}, \incluido{...} e
%% \substituido{...}{...} (e, para blocos inteiros, \paragraforemovido e
%% \reivindicacaoremovida). O mesmo arquivo gera as duas saídas.
%%
%% NÃO mexa no valor abaixo. Para gerar a cópia de comparação use
%%
%%     make comparacao        (ou ./gerar-copia-de-comparacao.sh)
%%
%% que produz os PDFs *-comparacao.pdf sem tocar neste arquivo. Deixar 'sim'
%% aqui faria os PDFs do pedido saírem marcados, contrariando o art. 57, I —
%% e por isso o `make verificar` recusa compilar assim.
%%%%%%%%%%%%%%%%%%%%%%%%%%%%%%%%%%%%%%%%%%%%%%%%%%%%%%%%%%%%%%%%%%%%%%%%

\providecommand{\CopiaDeComparacao}{nao}
\copiadecomparacao{\CopiaDeComparacao}

%%%%%%%%%%%%%%%%%%%%%%%%%%%%%%%%%%%%%%%%%%%%%%%%%%%%%%%%%%%%%%%%%%%%%%%%
%% 6. DADOS DO DEPÓSITO — NÃO SÃO IMPRESSOS EM NENHUM DOCUMENTO
%%
%% Guarde-os aqui para ter tudo em um lugar só, mas saiba que eles NÃO vão
%% para os PDFs: são preenchidos no formulário eletrônico de requerimento do
%% Peticionamento Eletrônico do INPI (art. 3º, I).
%%
%% O art. 21 proíbe, nos documentos do pedido, "timbres, logotipos, letreiros,
%% assinaturas ou rubricas, sinais ou indicações de qualquer natureza
%% estranhos à matéria do pedido". Nome de depositante, logotipo institucional,
%% endereço e assinatura, portanto, NÃO podem aparecer no relatório
%% descritivo, nas reivindicações, nos desenhos ou no resumo.
%%
%%   Depositante(s) ......... Empresa Brasileira de Pesquisa Agropecuária
%%   CNPJ ................... 00.348.003/0001-10
%%   Inventor(es) ........... (nomes completos, na ordem do requerimento)
%%   Procurador ............. (se houver; art. 59)
%%   Prioridade unionista ... (país, número e data do primeiro depósito; art. 12)
%%   Prioridade interna ..... (número e data do pedido anterior; art. 15)
%%   Período de graça ....... (forma, local e data da divulgação; art. 11)
%%   Código de serviço/GRU .. (natureza do pedido, paga até a apresentação)
%%   Figura que melhor representa o pedido ... Figura 1 (art. 37, § único)
%%   Listagem de sequências ... (se houver: código de controle devolvido pelo
%%                               Peticionamento ao anexar o XML no padrão ST.26;
%%                               Portaria INPI/PR nº 48/2022, arts. 4º e 6º)
%%%%%%%%%%%%%%%%%%%%%%%%%%%%%%%%%%%%%%%%%%%%%%%%%%%%%%%%%%%%%%%%%%%%%%%%
